\documentclass[UTF8,a4paper,11pt]{ctexart}

\usepackage{amsmath,amssymb}
\usepackage{geometry}
\geometry{margin=2.5cm}
\usepackage{booktabs}
\usepackage{hyperref}
\usepackage{enumitem}

\title{What are $A$, $\omega$, and $\phi$?}
\author{}
\date{}

\begin{document}
\maketitle

In

\[
x(t)=A\cos(\omega t+\phi)
\]

each symbol has a physical meaning:

\subsection{$A$ — Amplitude}

The maximum displacement from equilibrium.

\begin{itemize}
  \item If $A=5$ cm, the mass moves between $+5$ cm and $-5$ cm.
\end{itemize}

\bigskip\noindent\rule{\linewidth}{0.4pt}\bigskip

\subsection{$\omega$ — Angular Frequency}

Measures \textbf{how fast the oscillation occurs}.

\[
\omega = 2\pi f
\]

where $f$ is the frequency in Hz (cycles per second).

For a spring-mass system:

\[
\omega=\sqrt{\frac{k}{m}}
\]

Interpretation:

\begin{itemize}
  \item Larger spring constant $k$ $\rightarrow$ stiffer spring $\rightarrow$ faster oscillation.
  \item Larger mass $m$ $\rightarrow$ heavier object $\rightarrow$ slower oscillation.
\end{itemize}

Example:

\[
\omega=2\pi
\]

means the system completes 1 cycle every second.

\bigskip\noindent\rule{\linewidth}{0.4pt}\bigskip

\subsection{$\phi$ — Phase Angle (Initial Phase)}

Tells you \textbf{where the oscillation starts at $t=0$}.

For example:

\[
x(t)=A\cos(\omega t)
\]

At $t=0$,

\[
x(0)=A
\]

The mass starts at the top.

But if

\[
x(t)=A\cos\left(\omega t+\frac{\pi}{2}\right)
\]

then

\[
x(0)=0
\]

The mass starts at the center instead.

So $\phi$ shifts the wave left or right in time.

\bigskip\noindent\rule{\linewidth}{0.4pt}\bigskip

A useful visualization:

\begin{center}
\begin{tabular}{lll}
  \toprule
  Symbol & Meaning & Controls \\
  \midrule
  $A$ & Amplitude & How far it moves \\
  $\omega$ & Angular frequency & How fast it oscillates \\
  $\phi$ & Phase angle & Where it starts \\
  \bottomrule
\end{tabular}
\end{center}

Think of a Ferris wheel:

\begin{itemize}
  \item $A$ = radius of the wheel
  \item $\omega$ = how fast the wheel spins
  \item $\phi$ = where your seat is when the clock starts
\end{itemize}

That's why the solution of the spring equation $m\ddot{x}+kx=0$ is a sinusoid: the mass keeps exchanging energy between spring potential energy and kinetic energy, producing continuous oscillation.

\end{document}
